510=100
70 %

%TEXT%
Přidej sirku tak, aby platila rovnost.
\begin{center}
	\includegraphics[width=0.8\textwidth]{../data/Sirky/510=100.png}
\end{center}

%BAD%
\begin{enumerate}
\item Musí platit rovnost. Nelze vytvořit nerovnost.
\end{enumerate}

%HELP%
Musíte se dostat ze zajetých kolejí a podívat se na to trochu jinak

%SOLUTION%
\begin{center}
	\includegraphics[width=0.8\textwidth]{../data/Sirky/STO=100.png}
\end{center}

Pěkné, že? Nebo pořád nevidíte?\footnote{Číslo 510 je při zápisu sirkami stejné jako SIO. Stačí dát sirku nad I~a je z~něj T.}
